% ============================================================
\chapter{Multi-Cohort Integration and Biological Synthesis}
\label{chap:multi_cohort_integration}
% ============================================================

% ------------------------------------------------------------
\section{Rationale for Multi-Cohort Integration}
% ------------------------------------------------------------

The identification of early epigenetic alterations associated with the
Normal $\rightarrow$ Adjacent axis requires robustness across
independent studies. Rather than analysing each dataset in isolation,
the present chapter integrates GSE69914, GSE225845, and GSE287331
within a unified modelling framework.

The objective is not external validation between cohorts,
but the assessment of signal stability under controlled
multi-cohort pooling, where biological heterogeneity and
platform variability coexist within a single analytical setting.

% ------------------------------------------------------------
\section{Technical Challenges of Cohort Pooling}
% ------------------------------------------------------------

\subsection{Platform Heterogeneity}

The datasets originate from different Illumina platforms
(HumanMethylation450K and MethylationEPIC).
Differences in probe composition and genomic coverage
necessitate construction of a common CpG space prior to integration.

\subsection{Feature-Space Alignment}

A probe intersection strategy was applied to define a shared feature
space across cohorts.
All downstream modelling operates exclusively within this aligned set,
ensuring comparability across samples.

\subsection{Residual Cohort Effects}

Even after preprocessing standardisation,
cohort-specific variance structure,
clinical composition,
and sampling differences may induce latent effects.
The pooled modelling strategy therefore evaluates robustness
under heterogeneous but harmonised conditions.

% ------------------------------------------------------------
\section{Unified Preprocessing and Feature Selection}
% ------------------------------------------------------------

\subsection{Preprocessing Standardisation}

All datasets were processed using the identical pipeline
described in Chapter~\ref{chap:data_preprocessing},
including:

\begin{itemize}
\item technical probe filtering,
\item manifest-based annotation validation,
\item sex-chromosome exclusion,
\item $\beta \rightarrow M$ transformation.
\end{itemize}

This guarantees methodological uniformity prior to integration.

\subsection{Application of the Robust Feature Selection Framework}

The stability-driven feature selection procedure
introduced in Chapter~\ref{chap:feature_selection}
was applied within the pooled setting.
The framework incorporates:

\begin{itemize}
\item train-only variability filtering,
\item repeated stratified stability ranking,
\item effect-size prioritisation,
\item correlation-based redundancy pruning,
\item genomic diversification constraints.
\end{itemize}

The resulting CpG panel reflects loci that remain
statistically stable and directionally coherent
within a heterogeneous multi-cohort environment.

% ------------------------------------------------------------
\section{Pooled Multi-Cohort Modelling}
% ------------------------------------------------------------

\subsection{Joint Training Strategy}

Samples from all three cohorts were combined,
followed by stratified train–test splitting
performed on the aggregated dataset.
Model training was conducted on the pooled training set,
while evaluation was performed on the pooled test set.

This strategy does not constitute external validation,
but rather evaluates model robustness
under cohort heterogeneity.

\subsection{Performance Evaluation}

Model performance was assessed using:

\begin{itemize}
\item ROC–AUC,
\item balanced accuracy,
\item precision–recall metrics,
\item confusion matrix analysis.
\end{itemize}

Performance therefore reflects the capacity to capture
stable epigenetic drift patterns
shared across studies.

\subsection{Robustness Under Heterogeneity}

The pooled evaluation framework implicitly tests whether
the selected CpG panel encodes cohort-invariant
Normal–Adjacent differences,
rather than dataset-specific artefacts.

% ------------------------------------------------------------
\section{Gene-Level Aggregation and Functional Integration}
% ------------------------------------------------------------

\subsection{CpG-to-Gene Mapping}

Selected CpGs were mapped to gene symbols
using the Illumina manifest annotation.
Gene-level aggregation enables higher-order
biological interpretation.

\subsection{Gene Overlap and Functional Enrichment}

Over-representation analysis (ORA) was performed
to identify pathways enriched among the selected genes.
The background universe corresponds to the aligned CpG space.
False discovery rate correction was applied
to control multiple testing.

\subsection{Biological Coherence Across Cohorts}

Functional categories consistently enriched
in the pooled setting
indicate biological processes that remain stable
despite cohort heterogeneity.

% ------------------------------------------------------------
\section{Integrated Model of Early Epigenetic Drift}
% ------------------------------------------------------------

The combined evidence supports a model in which:

\begin{itemize}
\item early methylation alterations are focal rather than global,
\item directional effects persist across heterogeneous cohorts,
\item tumor samples display consolidation and increased variability,
\item adjacent tissue occupies an intermediate and unstable state.
\end{itemize}

Pooling cohorts strengthens statistical power
while revealing signals robust to study-specific variability.

% ------------------------------------------------------------
\section{Methodological Considerations}
% ------------------------------------------------------------

It is important to clarify that pooled evaluation:

\begin{itemize}
\item does not measure external transferability,
\item may partially absorb residual cohort effects,
\item reflects robustness within a harmonised multi-study framework.
\end{itemize}

Nevertheless, the stability-driven feature selection
and strict train-only procedures
mitigate information leakage
and enhance interpretability in HDLSS settings.

% ------------------------------------------------------------
\section{Synthesis of Multi-Cohort Evidence}
% ------------------------------------------------------------

In summary,
joint multi-cohort modelling confirms that
early epigenetic drift is subtle but detectable,
and that selected CpG panels retain discriminative ability
within a heterogeneous integrated framework.

This chapter consolidates the exploratory analyses,
preprocessing rationale,
and stability-based feature selection strategy,
providing a coherent biological synthesis
of early methylation alterations in breast tissue.